\section{Selección y adaptación del método}
La oficina de gestión utilizará dirección híbrida basada en PMBOK y construcción iterativa. La adaptación sigue el documento oficial PMI de 2021: elección de enfoque, ajuste organizacional, ajuste al proyecto y mejora continua. El método conserva controles contractuales e incorpora revisión frecuente de productos; la publicación fundamenta la adaptación y las decisiones siguientes pertenecen a Synaptix. \parencite[pp.~1--7]{sd6-pmi-adaptacion}

Proyecto coordina las líneas base; Funcional mantiene requisitos; Arquitectura controla diseño y compatibilidad; Datos, migración y conciliación; Seguridad, protección y riesgo; Calidad, evidencias; SRE, preparación y operación. El Jefe de Proyecto conserva dedicación exclusiva en implementación. La nominación y acreditación están en SD12/T-8 y las horas en T-15. La misma persona que cubra funciones distintas recibe una sola capacidad disponible. \parencite[cap.~19, pp.~33--34]{sd6-bt}

El expediente de adaptación identifica decisión, condición del caso, práctica elegida, responsable, producto, comprobación y documentos afectados. Su revisión mensual evalúa esperas, defectos, riesgo y carga del cliente. Un cambio de método que afecte el contrato se tramita antes de ejecutarse; una mejora interna conserva presupuesto de horas y puertas de recepción.

\section{Instancias de gobierno y comunicación}
La Tabla~\ref{tab:t9-comites} especifica las instancias del artículo 71. Proyecto prepara y distribuye el expediente; la secretaría designada levanta acta dentro de dos días hábiles, con acuerdos, responsables, fechas y referencias al registro afectado. \parencites[art.~71, pp.~37--38]{sd6-ba}[RT-19.09, p.~34]{sd6-bt}
\begingroup\setlength{\tabcolsep}{3pt}
\begin{longtable}{L{\dimexpr.20\textwidth-2\tabcolsep\relax}L{\dimexpr.16\textwidth-2\tabcolsep\relax}L{\dimexpr.32\textwidth-2\tabcolsep\relax}L{\dimexpr.32\textwidth-2\tabcolsep\relax}}
\caption{Participación y decisiones de gobierno}\label{tab:t9-comites}\\\hline
\EncabezadoTabla{Instancia} & \EncabezadoTabla{Cadencia} & \EncabezadoTabla{Participantes} & \EncabezadoTabla{Decisión y evidencia}\\\hline\endfirsthead
\hline\EncabezadoTabla{Instancia} & \EncabezadoTabla{Cadencia} & \EncabezadoTabla{Participantes} & \EncabezadoTabla{Decisión y evidencia}\\\hline\endhead
Ejecutivo & Mensual & Patrocinador CLIENTE, gerencia Synaptix y Administrador del Contrato. & Estrategia, escalamiento, cambios de alcance y riesgos mayores; acta y autorización escrita.\\\hline
Proyecto & Quincenal & Contraparte Técnica y Jefe de Proyecto. & Cronograma, entregables, desviaciones y decisiones operativas; acta y tablero.\\\hline
Arquitectura & Mensual & Arquitecto, Seguridad y referentes técnicos del CLIENTE. & Decisiones de arquitectura, deuda y riesgos técnicos; acta y registro de decisión.\\\hline
Operación & Mensual desde mes 13 & Líder de Operación, mesa de servicio y Contraparte Técnica. & Niveles de servicio, incidentes, problemas y mejora; reporte y acciones.\\\hline
Seguimiento & Semanal & Equipos de ambas partes. & Impedimentos y compromisos; registro de coordinación.\\\hline
\end{longtable}\FuenteTabla{Bases Administrativas, art.~71, pp.~37--38; Bases Técnicas Transversales, RT-19.09, p.~34.}\endgroup
La Contraparte designará referentes disponibles por proceso sin presumir un área informática del CLIENTE. Las decisiones de arquitectura se preparan con opciones comprensibles para esos referentes; el CLIENTE puede convocar apoyo competente. Proyecto da seguimiento a todo acuerdo hasta cerrar su evidencia, no sólo hasta emitir el acta.

Proyecto mantendrá el calendario de las cinco instancias y remitirá agenda, expediente y decisión solicitada con dos días hábiles de anticipación; una convocatoria extraordinaria registrará su urgencia. En el arranque se designarán titular y suplente de secretaría y, por cada rol participante, representación competente, disponibilidad y alcance de autoridad. La suplencia requiere designación escrita por la parte correspondiente y las acreditaciones exigidas para la función: no transfiere automáticamente facultades contractuales ni elimina la asistencia obligatoria. Proyecto comprobará la representación antes de la sesión y registrará ausencia, acción de cobertura y escalamiento al Administrador del Contrato. Synaptix incorporará reemplazo propio cuando la disponibilidad de los titulares no permita cubrir sus obligaciones, manteniendo dedicaciones y competencias de SD1/SD12.

La secretaría registrará ID y fecha de sesión, asistentes y representación, agenda, documentos y versiones examinados, acuerdos o desacuerdos, decisión y autoridad, responsable de cada acción, plazo, estado y enlace al riesgo, cambio o entregable. Distribuirá el borrador a participantes para revisión dentro del primer día hábil; incorporará las correcciones y publicará el acta dentro de los dos días hábiles siguientes a la sesión. Si persiste una discrepancia, la registrará expresamente junto con su responsable y plazo de resolución, sin aplazar la publicación ni atribuir aprobación por silencio. Una corrección posterior generará versión fechada con autor, motivo, cambio y comunicación a participantes; conservará el acta anterior y las autorizaciones firmadas, cuya modificación requiere a sus autoridades competentes. Proyecto revisará acciones en el tablero al cierre diario, advertirá antes del vencimiento y elevará los incumplimientos a la instancia con competencia. Sólo cerrará una acción con evidencia vinculada y confirmación del responsable de recepción.

La preparación, asistencia, revisión, corrección y seguimiento se imputan una sola vez a las dedicaciones de SD1 y a los paquetes de SD7/T-15. Los expedientes \texttt{1.1.4.m.q} de los meses 1--20 reservan 16 HH por quincena: Proyecto 8, Funcional 4 y Calidad 4; el control mensual de servicio \texttt{1.12.3.1.m} de 21--56 reserva 16 HH: Proyecto 4, Calidad 4 y Operación 8. Estas cifras son reservas de sus expedientes completos, no horas adicionales por reunión ni una reserva universal de todos los perfiles. El calendario de gobierno identificará, para cada sesión, preparación, asistencia y cierre por persona y paquete; Arquitectura, Seguridad, Operación y mesa enlazarán su intervención a sus paquetes y turnos correspondientes. Proyecto comprobará la suma con construcción, pruebas, implantación, mentoría, soporte y otras obligaciones, particularmente en 13--15 y 19--20. Las funciones compartidas usan una sola capacidad, y un agente de mesa no abandona su puesto sin relevo. Una insuficiencia obliga a incorporar cobertura propia y nivelar el registro antes de asignar el trabajo, sin duplicar las HH de un expediente o retrasar las cadencias contractuales.

El registro \texttt{INT-6} contiene interesado/rol, proceso, autoridad, impacto, disponibilidad, canal y responsable Synaptix. Administración valida saldos y contratos; torre, planes y avisos; portería, acceso; escuela, menores y retiros; varadero, maniobras y campañas. Funcional reúne las consultas antes de una sesión y registra la decisión o conflicto. Se evitan reuniones de validación durante maniobras, peak o congelamiento; una urgencia de seguridad se comunica por su canal de incidente. \parencite[numerales 13.1--13.3, pp.~28--30]{sd6-c07}

El espacio colaborativo mantiene documentos versionados, actas, entregables, riesgos y cambios, con acceso según perfil. El tablero se actualiza al cambiar un estado y al cierre diario. La comunicación semanal resume impedimentos; la quincenal, resultados y selección siguiente; la mensual, desempeño, hitos y riesgos. Proyecto es responsable de la publicación y Funcional de la coordinación con personas usuarias. \parencite[RT-19.05/06/08, pp.~33--34]{sd6-bt}

La programación y el control se calculan con programas reproducibles en Python 3 sobre los registros CSV UTF-8 de paquetes, red, calendario y asignaciones de SD7: se conservan versión del programa/entorno, entradas, resultados CPM/PERT, nivelación y validación de ciclos. El repositorio Git del CLIENTE versiona esa línea base, documentos y registros; el espacio colaborativo web autenticado publica el tablero HTML y permite descargar PDF/CSV. Proyecto publica al cambiar estado y al cierre diario, y verifica acceso de lectura del CLIENTE en cualquier momento; las aprobaciones se conservan como documentos firmados y no se sustituyen por un movimiento del tablero. Seguridad concede permisos mínimos, y cada actualización queda auditada por identidad e instante. La comprobación inicial reproduce una red, un corte de valor ganado y la vista del CLIENTE desde el mismo commit; respalda datos y procedimiento de publicación para restaurarlos. \parencites[Formulario T-7, Subdocumento 6, p.~58]{sd6-ba}[RT-19.05/08, pp.~33--34]{sd6-bt}

Funcional levantará requisitos mediante entrevistas por proceso, observación segura fuera de maniobras y talleres de escenarios con torre, portería, escuela, varadero y administración. Preparará antes de cada sesión reglas y datos de SD3/SD5; la salida será una ficha ITEM-6 con flujo normal, excepciones, criterios y decisiones abiertas. Los prototipos se revisarán con tareas concretas y permisos de datos adecuados. El referente funcional validará el resultado o registrará desacuerdo; Funcional consolidará conflictos y la Contraparte resolverá prioridades. La sesión se programa en los paquetes de especificación de SD7 y su resultado precede al inicio de construcción, sin sumar nuevamente las horas ya estimadas.

\section{Planificación, capacidad y configuración documental}
La selección de trabajo conserva los paquetes de SD7. Para cada resultado, Proyecto verifica criterio de aceptación, etapa, predecesoras, ejecutor, revisor y ventana. El esfuerzo asciende desde actividades, con contraste funcional UCP para los núcleos. Las reservas y esperas permanecen separadas; un porcentaje general no vuelve a sumar gestión o pruebas ya estimadas.

La red se comprueba sin ciclos; CPM identifica holgura total/libre y PERT considera tres duraciones. La nivelación considera calendario por persona, dedicación real, turnos y trabajo compartido. Proyecto no habilita un frente sin ejecución y revisión disponibles. La fecha de inicio usada para planificar es un supuesto de oferta hasta formalizarla; feriados, fechas náuticas y campañas efectivas se incorporan antes de aprobar la línea base.

La configuración documental conserva código de producto, versión, responsable, estado, fecha, solicitud que la originó y aprobación aplicable. T-12 mantiene íntegros sus 576 IDs y 27 complementos del caso; su estado cambia sólo tras evaluación por todas las cláusulas. La corrección metodológica no convierte un compromiso futuro en prueba ejecutada. Documentación vigente e histórica son distinguibles y cada liberación referencia versiones resueltas.

\section{Procedimiento de control de cambios}
El procedimiento \texttt{CC-6} cumple el artículo 72 y RT-19.03. Su registro contiene ID, solicitante, fecha, motivo, requisitos/paquetes afectados, clasificación, alternativas, impacto en alcance/plazo/costo/riesgo/servicio, decisión, responsables, versiones y comprobación. Los importes se conservan en el expediente económico autorizado. \parencites[art.~72, p.~38]{sd6-ba}[RT-19.03, p.~33]{sd6-bt}

La secuencia se ejecuta así:
\begin{enumerate}
\item Proyecto y Funcional registran la observación y clasifican defecto, precisión o modificación; mantienen su origen y recepción.
\item Arquitectura coordina impactos con Datos, Seguridad, Calidad, Integración y SRE; Proyecto comprueba red, capacidad, ventanas y efecto en hitos.
\item Se comprueba el límite acumulado de modificaciones del 20\,\% y las restricciones sobre objeto, naturaleza, garantías y cronograma del artículo 17.
\item El Comité Ejecutivo resuelve; ambas partes formalizan por escrito antes de ejecutar. Rechazo y aplazamiento se registran y comunican.
\item Proyecto publica las líneas base actualizadas, asigna trabajo autorizado y comunica la vigencia. Calidad verifica resultados y la Contraparte recibe el entregable aplicable.
\end{enumerate}
El cierre exige decisión escrita, versiones afectadas, ejecución trazable y comprobación. Aprobar una solicitud no significa aceptar el producto. Un defecto conserva su obligación de corrección; la reclasificación se justifica y no elimina requisitos de base.

La ejecución de cualquier cambio sin la aprobación previa del Comité Ejecutivo y el acuerdo escrito de ambas partes será de cargo y riesgo exclusivo de Synaptix y no dará derecho a pago adicional. Proyecto bloqueará su asignación como trabajo autorizado mientras falte cualquiera de esas dos decisiones; si se detecta ejecución no autorizada, registrará el incumplimiento, comunicará su impacto al Administrador del Contrato y coordinará contención y corrección seguras. La ejecución realizada no sustituye la aprobación ni altera retrospectivamente la línea base. El expediente conserva evidencia de decisión, acumulado calculado sobre el valor original del Contrato y comprobación de las materias inalterables, con acceso económico autorizado. \parencite[art.~72, numerales 1--5, p.~38]{sd6-ba}

\section{Medición de avance y recuperación}
El procedimiento \texttt{AV-6} vincula cada paquete con una cuenta de control y su presupuesto autorizado. Los hitos internos son especificación aceptada, construcción verificada y recepción del resultado; sus pesos se fijan desde el presupuesto de esas actividades antes de empezar y suman uno. Un hito gana su peso completo sólo con evidencia aprobada. No se remunera otra vez una prueba compartida ni se usa el número de commits como porcentaje de producto. Los servicios periódicos reconocen el trabajo efectivamente realizado en su período y se distinguen de productos discretos.

Con pesos $w_{ik}$ del paquete $i$, evidencias aceptadas $a_{ik}\in\{0,1\}$ y presupuesto $B_i$, el progreso verificable es $q_i=\sum_k w_{ik}a_{ik}$ y $EV=\sum_i B_iq_i$. $PV$ se obtiene de esos presupuestos con la programación aprobada al corte; $AC$, de registros reales conciliados. $SPI=EV/PV$ y $CPI=EV/AC$ usan unidades monetarias homogéneas; denominador cero impide calcular. El presupuesto, la regla de peso y los costos reales se documentarán en el ámbito contractual autorizado, sin publicarlos en la oferta técnica. \parencites[p.~1]{sd6-nasa}[RT-19.06/07, p.~34]{sd6-bt}[art.~50.2, p.~29]{sd6-ba}

El informe mensual incluye avance físico/financiero, cronograma, entregables, desviaciones, riesgos, incidencias y compromisos. Proyecto comprueba $q_i$ con Calidad y contrasta SPI con fechas y ruta crítica, porque SPI puede aproximarse a uno al terminar aunque exista atraso. Para el aviso RT-19.10 se usa $D_h=100\,\Delta t_h/T_h$, expresado en porcentaje, con atraso proyectado $\Delta t_h$ y duración aprobada $T_h$ del tramo hasta el hito; se publica la convención y además el atraso absoluto. Una desviación superior a 10\,\% exige aviso dentro de cinco días hábiles con plan de recuperación. Ninguna convención elimina el aviso oportuno de un incumplimiento contractual.

La recuperación declara causa, resultados en riesgo, alternativas, responsables y nueva proyección. Se verifica capacidad antes de paralelizar. Una nueva proyección no reemplaza la línea base sin autorización. El procedimiento RIESGO-8 de SD8 aplica ISO 31000:2018 y mantiene T-16 vivo, con causa, evento, impacto, evaluación inicial/residual, respuesta, responsable, plazo, disparador y prueba de eficacia. Proyecto lo revisa en cada Comité de Proyecto y ante aviso, incidente o cambio relevante; vincula cada acción a su paquete y registra consumo de reserva una sola vez. No se acepta como riesgo residual una excepción a una cláusula obligatoria. \parencites[RT-19.04/10, pp.~33--34]{sd6-bt}{sd6-iso31000}

\section{Adquisiciones y ventanas operacionales}
El procedimiento \texttt{ADQ-6} parte de la necesidad de SD7 y de T-11: especificación verificable, cantidad, variante, interfaces, fecha requerida y criterio de aceptación. Se comparan opciones por compatibilidad, soporte, licencia, seguridad, entrega, sustitución y dependencia. Se programa margen para abastecimiento y recepción. Las cifras de fabricante se sustentan en la variante real; una hipótesis de instalación no se presenta como catálogo.

Todo servicio externo declara alcance, acceso a datos, control, continuidad y condiciones de recepción; la autorización escrita y límite de subcontratación del artículo 73 se comprueban antes de contratar. La declaración de proveedores y alianzas específicas se mantiene en SD12. No se atribuye una carta de compromiso obtenida cuando se ofrece gestionarla. \parencite[art.~73, p.~38]{sd6-ba}

Ante aviso de marejadas se suspende inmediatamente toda actividad programada del proyecto, incluido trabajo aislado o remoto. Proyecto comunica el bloqueo, registra los paquetes y personas afectados y activa la respuesta PL-R01/R-01 de SD7/SD8. La reanudación exige comprobar cese de la condición, autorización y ventana segura; no basta cambiar de lugar el trabajo suspendido. Las liberaciones y capacitación respetan congelamiento total 15 de diciembre--15 de marzo, campañas de varadero abril--mayo y septiembre--noviembre y las 14 fechas náuticas. Los mantenimientos programados se anuncian al CLIENTE con al menos diez días hábiles; Proyecto conserva el aviso y SRE lo comprueba antes de liberar. La oficina aplica las reservas calculadas en SD7 y registra su consumo, sin duplicarlas ni trasladar hitos. \parencites[cap.~10, restricción 12, p.~26; numeral 13.2, p.~29; cap.~15, RT-10.05, p.~34]{sd6-c07}[art.~20, p.~15]{sd6-ba}[RT-10.05, p.~22]{sd6-bt}

\section{Comprobación del método y cierre}
Antes de usar el procedimiento, Proyecto y Calidad comprobarán una solicitud desde registro a decisión y una recepción desde paquete a prueba. El corte de avance debe ser reproducible por la Contraparte con acceso a sus evidencias. La prueba documental identifica responsables, versiones y resultado, sin suplir ensayos técnicos ni aceptación contractual.

El cierre por etapa exige productos, observaciones y aceptación de T-17/T-18. La marcha blanca se concilia diariamente y el cambio oficial de autoridad espera su acta. El expediente final identifica entregables, configuraciones, manuales, decisiones, riesgos residuales y transferencia a Operación; las lecciones se incorporan al registro de mejora. Las obligaciones de reversibilidad y salida se conservan en SD7, con responsables y calendario propios. \parencites[art.~17.3, pp.~12--13]{sd6-ba}[cap.~20, pp.~34--35]{sd6-bt}
